\documentclass[11pt]{article}
\usepackage[margin=1in]{geometry}
\usepackage{times}
\usepackage{hyperref}
\hypersetup{colorlinks=true, linkcolor=black, urlcolor=blue, citecolor=black}
\title{The Mach Effect Thruster and the Problem of Vibration Artifacts}
\author{Flyxion \\ \small Independent Researcher}
\date{}
\begin{document}
\maketitle

\begin{abstract}
The Woodward-Mach effect thruster differs from most reactionless-drive claims in this volume because its proponents publish a physical model, experimental methods, and a continuing response to criticism. The central claim is that a driven piezoelectric stack undergoes transient rest-mass changes that can be phased with internal motion to produce a net force. The most relevant independent analysis is the 2021 torsion-balance study by Monette, Kößling, and Tajmar at TU Dresden. Its abstract reports large second harmonics in the integrated strain-gauge signal, vibration coupling and amplification near system resonances, and measured forces that did not significantly change when drive conditions were varied, contrary to the model's prediction. The authors conclude that the observed effect in their experiment resulted from coupled vibrations on the torsion balance. James Woodward's later published response argues that vibration-only artifacts cannot produce a steady deflection consistent with momentum conservation, and reports a separate pendulum experiment in which the device and support moved together, which he interprets as real thrust. This response is a serious counterargument, but it does not reanalyze the Dresden signals or independently confirm their source. The present essay therefore treats the Dresden result as a specific null interpretation of its own torsion-balance measurements, not as proof that every Mach-effect device is impossible. The entropic-gravity framing remains brief because the apparatus tests a proposed inertial coupling, not a direct gravitational shielding effect.
\end{abstract}

\section{Introduction}

James Woodward's Mach-effect proposal begins with a controversial interpretation of inertia in general relativity. In the proposed model, changes in the internal energy of an accelerated body produce transient changes in its rest mass. If that variation is synchronized with a separate mechanical oscillation, the device is said to produce a small net force without expelling propellant. The apparatus most often discussed in the experimental literature uses a stack of piezoelectric elements driven electrically and mechanically constrained by end masses. The prediction is not simply that a vibrating object will move: the phase relation between the electrical drive and the mechanical motion is supposed to determine the direction and magnitude of a time-averaged force.

That claim is testable in a more specific way than many reactionless-drive proposals. The proponent literature includes experimental apparatus and proposed controls, including tests in which an arrangement intended to cancel the net effect is expected to suppress the reported thrust. Fearn and Woodward's 2013 arXiv paper describes one such null configuration: equal brass masses at either end of the stack should make the system push in both directions, and the authors report that the thrust signal disappears in that configuration \cite{FearnWoodward2013}. This is a result reported by the proponents, not an independent replication, but it shows that the claim has an experimentally stated control condition.

The more decisive question is what an independent group finds when it examines the force signal and the mechanical behavior of the balance itself. Monette, Kößling, and Tajmar at TU Dresden addressed that question in their 2021 paper on Mach-effect thrusters tested on torsion balances \cite{Monette2021}. Their conclusion is unusually direct: the force behavior they observed on their balance was the result of coupled vibrations, not evidence that the thruster generated the predicted Mach-effect force. Because this analysis targets the measurement apparatus and its vibration response, the essay is organized around its specific result and around the proponents' published reply.

\section{History of the Claim and Test Program}

The primary papers verified for this essay document two stages in the claim's experimental development rather than a complete history of every prototype. The 2013 paper by Fearn and Woodward presents an experiment and a null configuration in which the same capacitor stack is arranged to produce opposing effects at its two ends; its abstract reports that the thrust disappears in this configuration \cite{FearnWoodward2013}. This paper establishes what the proponents said they tested and what outcome they reported. It does not show that an independent laboratory confirmed the interpretation.

The 2021 Dresden study examined the electromechanical system on torsion balances and explicitly compared its measured forces with features predicted by Woodward's account \cite{Monette2021}. The Dresden authors note earlier claims and earlier experiments, but the history reconstructed here is limited to sources that could be checked directly. Woodward's 2021 IntechOpen chapter is a later proponent account that cites the Dresden study and describes a separate pendulum experiment intended to address vibration artifacts \cite{Woodward2021}. Thus, the available record includes an inventor-side account, a published independent analysis of a torsion-balance setup, and a subsequent proponent response reporting a different instrument. It does not amount to a consensus replication of the claimed effect.

\section{The Dresden Test Record}

The Dresden paper describes several linked vibration observations. Its abstract reports that the integrated strain-gauge signal from the prestressed piezoelectric stack contains large second harmonics, which the authors associate with vibration at system resonances and substantial nonlinearity. It further reports subharmonic vibration coupling in the balance-beam observations and amplified vibration near electromechanical resonances. The authors say that these effects can explain the transients present in the force trace. Finally, they report that changing the piezo actuator's driving conditions did not significantly change the measured force, although the Mach-effect theory predicts that the drive conditions should matter. Their abstract concludes that the observed effect was caused by coupled vibrations on the torsion balance \cite{Monette2021}.

This conclusion must be described at the level the paper supports. The Dresden group did not establish that every device built by Woodward or his collaborators must produce zero force under all conditions. It concluded that the signals measured in its torsion-balance investigation were explained by coupled vibrations and did not follow the expected dependence on drive conditions. The finding is consequential because it directly addresses whether the measured torsion-balance transients are reliable evidence for the claimed mechanism. It is not a universal proof against any possible Machian coupling.

The full paper makes the conclusion more specific than the abstract alone. Monette and colleagues report switching transients that reverse with device orientation, a small quasi-static component during fixed-frequency pulses, and larger responses during fine frequency sweeps. The fixed-frequency component was two orders of magnitude below Woodward's claimed effect. The larger sweep responses also appeared with the device at 90 degrees, which the authors say disagreed with the predicted uniaxial thrust direction. They further report that measured force did not depend strongly on drive frequency, voltage, or second-harmonic content as the model predicted. Low-frequency beam vibration was observed in all tested orientations, near resonances, and during changing drive conditions, particularly at switch-on. Their conclusion states that the effects observed with Woodward's MEGA drive on a torsion balance could be explained by thermal and vibrational artifacts using Newtonian mechanics \cite{Monette2021}.

The distinction between vibration and thrust is technically important. A driven piezo stack is designed to vibrate, and its mount, beam, strain gauge, wiring, and support form a coupled mechanical system. A sensor can register force-like transients when a component vibrates against a compliant or resonant structure even if the whole apparatus has not acquired a sustained external impulse. A vibration at a structural resonance may be amplified; nonlinear response can generate harmonics that do not track the idealized drive in a simple way. The Dresden paper also discusses a possible center-of-mass offset on the balance, thermal expansion near resonances, and loosening or stick-slip at screw connections as mechanisms that can create offsets or switching transients. The authors do not claim that each proposed artifact was independently proven to account for every signal. Their stated conclusion is that thermal and vibration effects can explain the observed behavior in their balance. In that setting, a force trace is not self-interpreting. The experimental question is whether the signal follows the device's predicted phase and control dependence, and whether the entire thruster-plus-support system behaves as a body receiving net external momentum.

The Dresden authors' point about drive-condition dependence strengthens their interpretation. If the proposed effect depends on a particular phase and operating condition, a genuine signal should respond to deliberate changes in those variables in a way specified in advance by the model. The reported near-invariance of the measured force under changes that should matter is therefore an anomaly for the proposed explanation, even if the apparatus shows repeatable mechanical motion. It shifts the evidentiary burden from showing that a balance moves to showing that its motion has the signature expected from the Mach-effect model rather than from vibration coupling.

\section{The Strongest Objection: Woodward's Published Response}

The strongest response from the proponents is James Woodward's published chapter, ``Gravity and Inertia in General Relativity,'' which directly cites the Dresden paper and addresses vibration artifacts as the leading alternative explanation for Mach-effect force claims \cite{Woodward2021}. Woodward argues that vibration alone cannot sustain a steady sensor displacement in a system that conserves momentum. In his account, a vibration artifact should initially move the device and its support in opposite directions, so that the center of mass remains unchanged, whereas a real net force should cause the device and support to move together and leave a continuing offset while power is applied.

Woodward reports a separate pendulum experiment rather than a reanalysis of the Dresden torsion-balance data. He describes mounting a Mach-effect gravity-assist device on a pendulum platform and measuring both the position of the device and the support structure. He states that the two position traces moved together, which he interprets as evidence of a real force rather than a vibration artifact. He also argues that the observed platform displacement was much larger than the vibration amplitudes measured on the device, making a stick-slip account insufficient; he reports a dummy-capacitor control that did not produce similar pendulum motion. In the illustrated run he calculates about 250 micronewtons from the pendulum geometry, while his conclusions summarize the force as about 200 micronewtons \cite{Woodward2021}. These quantities are reported by Woodward and should be read as his account of his own experiment, not as an independently established measurement.

This reply deserves fair consideration. It articulates a discriminating experimental signature, identifies a specific mechanical alternative, and reports controls intended to separate the two interpretations. It is more substantial than simply asserting that a vibration explanation is impossible. The central limitation is that it concerns a different apparatus and a separate experiment. It does not show that the harmonics, resonance coupling, and force-trace transients observed by Monette and colleagues were not caused by the coupled vibrations they identified. Nor does a force measured and interpreted by the proponents independently validate the Dresden apparatus or overturn its reported null interpretation.

\section{What the Disagreement Does and Does Not Establish}

The two records answer related but distinct experimental questions. The Dresden group reports that its torsion-balance signal was dominated by coupled vibration effects and did not respond as the model predicted when the drive conditions changed. Woodward responds that a vibration-only model cannot explain a steady deflection in a different pendulum setup where the device and its support moved together, and he presents that behavior as evidence for a real force. The disagreement should not be reduced to the claim that Dresden ignored the proponents' argument: Woodward's chapter cites their torsion-balance paper and addresses vibration artifacts directly. But his chapter does not reproduce their data or demonstrate that their measured transients were thrust.

The strongest remaining defense of the Mach-effect claim is that the Dresden group's balance may be particularly susceptible to vibration coupling, while another support design and a different sensor geometry can separate a real net force from internal vibration. The proponents' pendulum test is intended to answer exactly that objection. The strongest answer from the Dresden analysis is that the measured signal in its own setup has identified vibration signatures, including resonance-associated harmonics and force transients, and does not display the drive dependence the model predicts. Neither point can be substituted for an independent replication of the other group's procedure.

A more decisive comparison would use the same thruster under both measurement approaches, with the predicted phase and drive variations specified before the run, simultaneous tracking of the device and support, direct monitoring of structural vibration, and positive and null configurations. Ideally, the data and analysis would be reviewed by a team not invested in either outcome. Such a test should not be described as already completed unless a report documents it. The published record reviewed here instead contains a clearly stated independent artifact interpretation and a proponent-authored response from a distinct experiment.

\section{The Entropic Gravity Framing}

The Mach-effect thruster invokes an interpretation of inertia as sensitive to distant matter and gravitational coupling. The entropic-gravity framing used in this series is therefore relevant only in a limited sense. An emergent account of gravitation may be compatible with global dependence in principle, but it does not establish the magnitude, phase dependence, or laboratory thrust predicted by Woodward's model. Those are questions for the model and the experiment. In this case the immediate issue is whether the reported force signals are genuine net external forces or vibration artifacts in the measuring system.

\section{Conclusion}

The Dresden group's 2021 torsion-balance paper concluded that the effects observed in its experiment resulted from coupled vibrations: the strain-gauge signal showed resonance-related second harmonics, balance-beam oscillations showed vibration coupling and amplification, and measured forces did not significantly change under drive variations that the theory predicts should matter. Woodward's published response argues that vibration artifacts cannot sustain a steady deflection and reports a separate pendulum experiment in which the device and support moved together, which he interprets as a Mach-effect force. That is a substantive response, but it is a proponent's interpretation of another apparatus and does not directly explain away the Dresden measurements. This paper does not prove that every Mach-effect thruster is incapable of producing force, nor does it independently validate Woodward's pendulum result. It concludes that the Dresden signal is not persuasive evidence for the claimed effect under the conditions they tested, while the broader claim remains dependent on independent tests that distinguish the specific vibration signatures from a predicted net thrust.

\begin{thebibliography}{99}

\bibitem{FearnWoodward2013} H. Fearn and J. F. Woodward, ``Experimental Null Test of a Mach Effect Thruster,'' arXiv:1301.6178, 2013. https://arxiv.org/abs/1301.6178.

\bibitem{Monette2021} M. Monette, M. K\"o{\ss}ling, and M. Tajmar, ``Experimental Investigation of Mach-Effect Thrusters on Torsion Balances,'' \emph{Acta Astronautica} 182 (2021), 531--546. https://doi.org/10.1016/j.actaastro.2021.02.037.

\bibitem{Woodward2021} J. F. Woodward, ``Gravity and Inertia in General Relativity,'' IntechOpen, 2021. https://doi.org/10.5772/intechopen.99760.

\end{thebibliography}

\end{document}
